\documentclass[12pt]{article}
\usepackage[english]{babel}
\usepackage{geometry}
\geometry{a4paper, margin=1in}
\usepackage{setspace}
\doublespacing
\usepackage{fontspec}
\setmainfont{Times New Roman}
\usepackage{indentfirst}
\usepackage[notes]{biblatex-chicago}

\title{The Superiority of the Classical Hippocratic Oath: A Defense Against Modern Critiques}
\author{JC Vaught}
\date{\today}

\begin{document}

\maketitle

\section*{Introduction}

The Classical Hippocratic Oath, composed in ancient Greece, has endured for millennia as a moral foundation for medicine. Its principles—reverence for life, fidelity to mentors, and sacred duty—stand in stark contrast to the Modern Hippocratic Oath, which emphasizes adaptability, patient autonomy, and social responsibility.\footnote{Hippocrates, \textit{Hippocratic Oath: Classical Version}, trans. Ludwig Edelstein (Chicago: University of Chicago Press, 1943).} Critics argue that the classical oath is outdated, rigid, and incompatible with contemporary values. However, its unwavering ethical boundaries, emphasis on tradition, and spiritual grounding offer indispensable safeguards against the moral relativism and systemic flaws of modern healthcare. This essay defends the classical oath’s superiority by addressing and refuting common critiques, demonstrating that its perceived limitations are, in fact, its greatest strengths.

\section*{Unambiguous Moral Boundaries vs. Charges of Inflexibility}

The classical oath establishes clear, non-negotiable prohibitions: “I will neither give a deadly drug to anybody who asked for it, nor will I make a suggestion to this effect. Similarly, I will not give to a woman an abortive remedy.”\footnote{Hippocrates, \textit{Classical Version}.} These absolutes prioritize the sanctity of life, shielding physicians from ethical ambiguity. By categorically rejecting euthanasia and abortion, the oath ensures that medical decisions align with the fundamental duty to preserve life, even in morally fraught situations.

Critics of the classical oath argue that its rigidity ignores patient autonomy and modern societal values. The Modern Hippocratic Oath, by contrast, permits euthanasia in cases of terminal suffering and allows abortion under specific conditions, framing these acts as compassionate responses to individual circumstances.\footnote{Louis Lasagna, \textit{Hippocratic Oath: Modern Version} (Boston: Harvard Medical School, 1964).} This flexibility, opponents claim, respects bodily autonomy and avoids harm caused by dogmatic prohibitions.

However, the classical oath’s inflexibility is not a weakness but a bulwark against moral compromise. History demonstrates that ethical systems permitting exceptions—even for compassionate reasons—risk normalizing practices that devalue life. For instance, legalized euthanasia, initially framed as a “mercy” for the terminally ill, has expanded in some countries to include non-terminal psychiatric patients, raising concerns about coercion.\footnote{Hippocrates, \textit{Classical Version}.} The classical oath’s prohibitions protect physicians from external pressures to instrumentalize death, ensuring they act as healers rather than arbiters of life and death.

\section*{Mentorship and Tradition vs. Charges of Exclusion}

The classical oath binds physicians to a closed network of mentorship, requiring them to teach only those who “have signed the covenant and have taken an oath according to the medical law.”\footnote{Ibid.} This system ensures apprentices inherit both technical skills and the moral ethos of their predecessors, fostering accountability. 

Modern critics condemn this exclusivity as elitist, arguing that the Modern Oath’s inclusive approach—“gladly shar[ing] such knowledge as is mine with those who are to follow”—better aligns with democratic ideals.\footnote{Lasagna, \textit{Modern Version}.} Yet the classical model’s exclusivity is not gatekeeping but a safeguard against the commodification of medical education. By restricting teaching to those who swear the oath, it filters out individuals motivated by profit rather than duty.\footnote{Hippocrates, \textit{Classical Version}.}

\section*{Sacred Duty vs. Charges of Irrelevance in a Secular World}

Invoking divine witnesses, the classical oath imbues medicine with spiritual significance: “I swear by Apollo Physician and Asclepius... in purity and holiness I will guard my life and my art.”\footnote{Ibid.} Secular critics dismiss this religiosity as outdated, asserting that ethics should derive from humanistic principles.\footnote{Lasagna, \textit{Modern Version}.} Yet the oath’s spiritual language underscores the gravity of medical practice, deterring hubris in an era of technological arrogance. The vow to avoid “sexual relations” with patients, for example, is a pledge of moral purity, not merely a rule.\footnote{Hippocrates, \textit{Classical Version}.}

\section*{Conclusion}

The Classical Hippocratic Oath’s critics misunderstand its purpose. It is not a relic but a manifesto for preserving medicine’s soul. Its unambiguous boundaries prevent ethical slippage, its mentorship model sustains moral continuity, and its sacred ethos guards against dehumanizing commodification. For a profession entrusted with humanity’s most vulnerable, the classical oath remains not just relevant but essential.

\end{document}